\chap[intro] Introduction

Passkeys have become a popular means of authentication on the Internet. Along with more conventional means of passkeys, users are also adopting hardware security keys as a safer way to store their credentials. Hardware security keys may provide an additional layer of security compared to platform authenticators. They offer stronger isolation of private keys and ensure their non-exportability, thereby reducing the risk of credential compromise even if the host device is not trusted.

\nl One such project is the LionKey, an open-source hardware security key developed at CTU FEE by Ing.~Martin Endler. The device is compliant with FIDO2 standards; however, it currently supports only CTAP over USB (CTAPHID). This means that the key cannot be used with devices that do not support USB-based communication, notably mobile phones.

\sec[objective] Thesis Objective

The objective of this thesis is to extend the communication capabilities of the LionKey hardware security key by implementing an NFC-based communication interface on top of the existing CTAPHID solution. The goal is to enable interoperability with devices that lack USB support, particularly mobile platforms.

\nl The NFC interface is realised using the STEVAL-D25R3916B board, which is interfaced with the ST Nucleo-H533RE microcontroller in accordance with the original specifications of the LionKey project.

\sec[structure] Thesis Structure

The thesis is organized into six main parts:
\nl The Introduction outlines the context of the work and formulates the objectives of the thesis.
\nl Chapter~\ref[passkeys-theory] presents the theoretical background of passkeys and motivates their use in modern authentication systems.
\nl Chapter~\ref[nfc-theory] introduces the fundamental principles of NFC technology, focusing on aspects relevant to the proposed solution.
\nl Chapter~\ref[implementation] details the design, hardware used, and implementation of the proposed solution.
\nl Chapter~\ref[evaluation] describes the results achieved by this implementation and its applicability to practice.
\nl Chapter~\ref[conclusion] summarizes the results achieved and discusses the general outcome of the thesis.
